% based on a template made by the university of cologne
% http://www.mi.uni-koeln.de/wp-MIEDV/wp-content/uploads/2016/07/LaTeX-Vorlage.zip - 2023-11-02
\documentclass[12pt,a4paper]{scrartcl}

\addtokomafont{sectioning}{\rmfamily}
\usepackage[ngerman]{babel}% deutsches Sprachpaket wird geladen
\usepackage[T1]{fontenc} % westeuropäische Codierung wird verlangt
\usepackage[utf8]{inputenc}% Umlaute werden erlaubt
\usepackage[usenames]{color} % Erlaubt die Benutzung der namen im Farbpaket und deren Änderung
\usepackage{amsmath} % Erweiterung für den Mathe-Satz
\usepackage{amssymb} % alle Zeichen aus msam und msmb werden dargestellt
\usepackage{graphicx} % Graphiken und Bilder können eingebunden werden
%\usepackage{multirow} % erlaubt in einer Spalte einer Tabelle die Felder in mehreren Zeilen zusammenzufassen
\usepackage{enumerate} % erlaubt Nummerierungen
\usepackage{xurl} % Dient zur Auszeichnung von URLs; setzt die Adresse in Schreibmaschinenschrift.
\usepackage[center]{caption}  % Bildunterschrift wird zentriert
%\usepackage{subfigure} % mehrere Bilder können in einer fugure-Umgebung verwendet werden
%\usepackage{longtable} % Diese Umgebung ist ähnlich definiert wie die tabular-Umgebung, erlaubt jedoch mehrseitige Tabellen.
%\usepackage{paralist} % Modifikation der bereits bestehenden Listenumgebungen
\usepackage{lmodern}% Für die Schrift
\usepackage[hidelinks]{hyperref} % Links und Verweise werden innerhalb von PDF Dokumenten erzeugt
%\usepackage{wrapfig} % Das Paket ermöglicht es von Schrift umflossene Bilder und Tabellen einzufügen.
\usepackage{latexsym} % LaTeX-Symbole werden geladen
\usepackage{tikz} % Erlaubt es mit tikz zu zeichnen
\usepackage{tabularx} % Erlaubt Tabellen
\usepackage{algorithm} % Erlaubt Pseudocode
\usepackage{color} % Farbpaket wird geladen
%\usepackage{stmaryrd} % St Mary Road Symbole werden geladen
\usepackage{physics}
\usepackage{mhchem} % Chemie: \ce & \pu

\numberwithin{equation}{section} % Nummerierungen der Gleichungen, die durch equation erstellt werden, sind gebunden an die section
\newcommand{\HRule}{\rule{\linewidth}{0.7mm}}
\newcommand{\pu}[1]{\ensuremath{\mathrm{#1}}}

% disable commands
\renewcommand{\[}{} % math block start
\renewcommand{\]}{\noindent} % math block end
\newcommand{\tightlist}{} % created in enumerations

\hypersetup{
  pdftitle={B3.4},
  pdfcreator={LaTeX via pandoc}}

\setcounter{secnumdepth}{6}
\setcounter{tocdepth}{6}

\begin{document}
\begin{titlepage}
	\pagestyle{empty}

	\begin{center}

	\textsc{\LARGE Universität zu Köln }\\ [0.4cm]
	\textsc{Mathematisch-Naturwissenschaftliche Fakultät} \\[1.5cm]

	\includegraphics[width=0.45\textwidth]{../media/uni.jpg}\\[1.5cm]  % Uni-Logo wird geladen

	\textsc{\Large Praktikum~B}\\[2mm]
	\textsc{}\\[10mm]
	\HRule \\[0.4cm]

		{	\Huge \bfseries B3.4}\\[0.4cm]
			{	\huge \bfseries Positronen--Emissions--Tomographie}\\[0.3cm]
	
	\HRule \\[3cm]

 	\begin{center}
		\textsc{\Large Catherine~Tran } \\[3pt]
		\textsc{\Large Carlo~Kleefisch } \\[3pt]
		\textsc{\Large Oliver~Filla } \\[3pt]
	\end{center}
	\end{center}
\end{titlepage}

\newpage
\tableofcontents
\newpage

\hypertarget{einleitung}{%
\section{Einleitung}\label{einleitung}}

Die Positronen--Emissions--Tomographie (PET) ist ein
nuklearmedizinisches bildgebendes Verfahren, bei denen radioaktive
Materialien als \emph{Tracer} verabreicht werden, die dann im Körper des
Patienten gemessen werden. Dadurch können Bilder von z.B. Krebszellen
erzeugt werden.

In diesem Versuch wird eine radioaktive Quelle in einer Probe mittels
der PET lokalisiert. Weiterhin wird die Ortsauflösung sowie die
Winkelabhängigkeit untersucht.

\clearpage
\hypertarget{theoretische-grundlagen}{%
\section{Theoretische Grundlagen}\label{theoretische-grundlagen}}

\hypertarget{pet}{%
\subsection{PET}\label{pet}}

Die Positronen--Emissions--Tomographie ist ein Verfahren, das die
Möglichkeit bietet, Veränderungen von Stoffwechselprozessen sowie
physiologische Aktivitäten sichtbar zu machen. In der Medizin wird
dieses Verfahren hauptsächlich verwendet, um Krebszellen im Körper zu
detektieren.

Hierbei injiziert man dem Patienten radioaktive Markerstoffe
(\emph{Tracer}). Diese werden unterschiedlich von den Zellen aufgenommen
und verarbeitet, abhängig vom lokalen Stoffwechsel. Damit der Körper
nicht zu viel radioaktive Strahlungen ausgesetzt wird, muss der Tracer
so kurzlebig wie möglich sein und dennoch Positronen abgeben. Deswegen
eignen sich z.B. \(\ce{^{11}C}\), \(\ce{^{13}N}\), \(\ce{^{15}O}\),
\(\ce{^{18}F}\) besonders. Unterschiedliche Tracer werden für
unterschiedliche Körperteile verwendet.

Wenn der Tracer im Gewebe ist, zerfällt er über den \(\beta^+\)--Zerfall
und emittiert dabei Positronen. Innerhalb weniger Mikrometer
rekombiniert ein Positron mit einem Elektron und es kommt zur
\emph{Paarvernichtung}. Hierbei entstehen zwei Photonen mit jeweils
einer Energie von \(\pu{511 keV}\) ein einem Winkel von \(180^\circ\)
zueinander. \([6]\)

Um die Aktivität quantitativ zu bestimmen, werden Signale von den
ausgestrahlten Photonen registriert. Wenn zwei davon koinzident sind,
hat dieser Prozess mit großer Wahrscheinlichkeit auf der
Verbindungslinie\footnote{Englisch: \emph{line of response} (\emph{LOR})}
stattgefunden. Aus vielen solcher koinzidenten Signale lässt sich die
Aktivität mit einem Computerprogramm zu einem zweidimensionalen und
später dreidimensionalen Bild rekonstruieren.

Es kann bei der Messung häufig zu falschen Koinzidenzen kommen. Entweder
sind das nur zufällige Signale, die zeitgleich aufgenommen werden, oder
das Photon wurde schon auf den Weg zum Detektor gestreut. Diese falsche
Koinzidenzen möchte man möglichst gering halten, da sie nicht zur
Bildqualität tragen und Nachbearbeitung fordern.

Neben der PET gibt es noch andere bildgebende Verfahren wie
Röntgentomographie, Computertomographie (CT), Magnetresonanztomographie
(MRT) und Ultraschall. Die Strahlenbelastung ist bei jeweiligem
Verfahren unterschiedlich.

Bei Ultraschall werden nur unschädliche Schallwellen verwendet, MRT ist
ebenso frei von Strahlenbelastung. Im Vergleich zu Röntgentomographie
und CT ist PET weniger schädlich, da die Stoffe sehr schnell zerfallen.
Eine Röntgenaufnahme am Brustkorb hat eine Äquivalentdosis von
\(\pu{0.2 mSv}\) und eine CT am Brustkorb hat ca \(\pu{8 mSv}\) \([8]\).

Laut Bundesamt für Strahlenschutz beträgt der Grenzwert für eine
einzelne Person \(\pu{1 mSv}\) pro Jahr, so lange man nicht beruflich
mit Strahlung zu tun hat.

Für die Medizin ist PET ein hervorragendes Werkzeug, besonders in der
Onkologie. Das Verfahren kann Veränderung im Stoffwechsel erkennen, eher
ein anderes Verfahren es nachweisen kann, so können Krebstumoren früher
erkannt und behandelt werden.

\hypertarget{betazerfall}{%
\subsection{\texorpdfstring{\(\beta\)--Zerfall}{\textbackslash beta--Zerfall}}\label{betazerfall}}

Der \(\beta\)--Zerfall ist eine der drei Arten radioaktiven Zerfalls.
Wie beim \(\alpha\)--Zerfall wird dabei ein chemisches Element in ein
anderes umgewandelt. Im Unterschied zum \(\alpha\)--Zerfall werden
hierbei keine Nukleonen abgesondert, sondern ein Nukleon in ein anderes
umgewandelt.

Es wird zwischen drei Arten des \(\beta\)--Zerfalls unterschieden. Es
gibt \(\beta^+\)--Zerfall und \(\beta^-\)--Zerfall, weiterhin wird der
Elektroneneinfang dazugezählt.

Radioaktiver Zerfall geschieht, wenn der Tochterkern eine höhere
Bindungsenergie als der Mutterkern erhält. Hierfür ist insbesondere der
Symmetrieterm der Weizsäcker Massenformel \((??)\) interessant. Daran
kann man erklären, dass \(\beta^-\)--Zerfall dann geschehen kann, wenn
der Kern mehr Neutronen als Protonen hat.

Umgekehrt kann es zu \(\beta^+\)--Zerfall kommen, wenn der Kern mehr
Protonen als Neutronen hat. Ebenso ist Elektroneneinfang möglich.
Welcher dieser Prozesse stattfindet hängt von der Energiedifferenz
zwischen Mutterkern und Tochterkern ab: Wenn die Energiedifferenz
kleiner als \(\Delta E_\mathrm{min}=\pu{1.022 MeV}\) beträgt, dann
reicht die Energie nur für Elektroneneinfang. Bei einer Energie von mehr
als \(\Delta E_\mathrm{min}\) ist auch \(\beta^+\)--Zerfall möglich.
\(\Delta E_\mathrm{min}\) entspricht der doppelten Ruheenergie von
Elektronen.

Beim \(\beta^+\)--Zerfall wird ein Proton \(p\) in ein Neutron \(n\)
umgewandelt, dabei entstehen ein Positron \(e^+\) sowie ein
Elektronenneutrino \(\nu_e\). Hierbei wird ein Mutterkern
\(\ce{^A_Z X}\) in einen Tochterkern \(\ce{^A_{Z-1} Y^-}\) umgewandelt.
Beim \(\beta^-\)--Zerfall wird ein Neutron \(n\) in ein Proton \(p\)
umgewandelt, dabei entstehen ein Positron \(e^+\) sowie ein
Anti--Elektronenneutrino \(\bar{\nu}_e\). Hierbei wird ein Mutterkern
\(\ce{^A_Z X}\) in einen Tochterkern \(\ce{^A_{Z+1} Y^+}\) umgewandelt.

Beim Elektroneneinfang (ec) wird ähnlich wie beim \(\beta^+\)--Zerfall
ein Proton \(p\) in ein Neutron \(n\) umgewandelt. Allerdings wird dazu
ein Elektron \(e^-\) verwendet, dass aus der \(K\)--Schale des Atoms
eingefangen wurde. Daher wird nur ein Elektronenneutrino \(\nu_e\)
erzeugt. Hierbei wird ein Mutterkern \(\ce{^A_Z X}\) in einen
Tochterkern \(\ce{^A_{Z-1} Y}\) umgewandelt.

\[
\begin{eqnarray}
    \beta^+:\qquad\,\,\qquad
        \ce{p} &\ce{->}& \quad \ce{n} \quad\, + \ce{e^+} + \ce{\nu_e} \\
        \ce{^A_Z X} &\ce{->}& \ce{^A_{Z-1} Y^-} + \ce{e^+} + \ce{\nu_e} \\
    \beta^-:\qquad\,\,\qquad
        \ce{n} &\ce{->}& \quad \ce{p} \quad\, + \ce{e^-} + \ce{\bar{\nu}_e} \\
        \ce{^A_Z X} &\ce{->}& \ce{^A_{Z+1} Y^+} + \ce{e^-} + \ce{\bar{\nu}_e} \\
    \text{ec}:\qquad
        \ce{p + e^+} &\ce{->}& \quad \ce{n} \ \ \, + \ce{\nu_e} \\
        \ce{^A_Z X} &\ce{->}& \ce{^A_{Z-1} Y} + \ce{\nu_e}
\end{eqnarray}
\]

\hypertarget{paarvernichtung}{%
\subsection{Paarvernichtung}\label{paarvernichtung}}

Der Begriff der Paarvernichtung beschreibt die Zerstrahlung von einem
Teilchen und seinem Antiteilchen, bei der die Teilchen in
\(\gamma\)--Strahlung umgewandelt werden. Die Energie der Strahlung ist
die Summe der kinetischen Energie dieser Teilchen sowie ihrer
Ruhemassen. Die Paarerzeugung ist der dazu inverse Effekt. \([5]\)

Bei der Vernichtung eines Positrons \(e^+\) mit einem Elektron \(e^-\)
entstehen normalerweise zwei \(\gamma\)--Quanten in einem Winkel
\(\theta\). Dieser hängt von der transversalen Impulskomponente \(p_T\)
ab. Weiterhin sind die Ruhemasse \(m_e\) der Teilchen und die
Lichtgeschwindigkeit \(c\) relevant. \([6]\)

\[
\begin{eqnarray}
    \tan(\theta) &=& \frac{p_T}{m_ec}
\end{eqnarray}
\]

Sind die beiden Teilchen in Ruhe zueinander, so ist der Abstahlwinkel
der \(\gamma\)--Quanten \(\pu{180^\circ}\). Haben sie einen relativen
Impuls, so wird der Winkel kleiner. Sind die Teilchen jedoch zu schnell,
dann ist der Wirkungsquerschnitt sehr klein und es ist sehr
unwahrscheinlich, dass Paarvernichtung stattfindet.

\hypertarget{gammastrahlung}{%
\subsection{\texorpdfstring{\(\gamma\)--Strahlung}{\textbackslash gamma--Strahlung}}\label{gammastrahlung}}

Gammastrahlung ist im engeren Sinne eine besonders durchdringende
elektromagnetische Strahlung, die bei radioaktivem Zerfall auftritt. Sie
wirkt durch drei Prozesse auf Materie. Bei einer Energie von bis zu
\(E_\gamma = \pu{100 keV}\) überwiegt der Photoeffekt, weiterhin gibt es
den Comptoneffekt und Paarbildung.

\hypertarget{photoeffekt}{%
\subsubsection{Photoeffekt}\label{photoeffekt}}

Beim Photoeffekt überträgt ein Photon seine gesamte Energie auf ein
Hüllenelektron, welches daraufhin das Atom verlässt. Er wirkt vor allem
bei elektromagnetischer Strahlung mit einer Energie von bis zu
\(E_\gamma = \pu{100 keV}\).

Die kinetische Energie \(E_\mathrm{kin}\) des Elektrons hängt von der
Energie des Photons \(h\nu\) sowie der Austrittsarbeit \(W_A\), ab.
Letztere ist notwendig, um das Elektron aus dem Atom zu lösen.

\[
\begin{eqnarray}
    E_\mathrm{kin} &=& h\nu - W_A
\end{eqnarray}
\]

Der Wirkungsquerschnitt \(\sigma_\mathrm{ph}\) wird durch die
Ordnungszahl \(Z\), die reduzierte Photonenenergie \(\epsilon\), die
Sommerfeld'sche Feinstrukturkonstante \(\alpha\) und den
Thomson--Wirkungsquerschnitt \(\sigma_\mathrm{Th}^e\) bestimmt. \([7]\)

\[
\begin{eqnarray}
    \sigma_\mathrm{ph}
        &=& \sqrt{\frac{32}{\epsilon^7}}\alpha^4 Z^5
            \sigma_\mathrm{Th}^e \pu{\frac{cm}{Atom}} \\
    \sigma_\mathrm{ph}
        &\propto& \frac{Z^5}{E_\gamma^{\frac{7}{2}}} \\
    \epsilon &=& \frac{E_\gamma}{m_ec^2} \\
    \sigma_\mathrm{Th}^e &=& \frac{8}{3} \pi r_e^2
\end{eqnarray}
\]

Hierbei sind weiterhin die Elektronenmasse \(m_e\) und der Radius
\(r_e\) eines Elektrons relevant.

\hypertarget{comptoneffekt}{%
\subsubsection{Comptoneffekt}\label{comptoneffekt}}

Der Comptoneffekt tritt vor allem bei \(\gamma\)--Strahlung mit
mittleren Energien von ca. \(\pu{100 keV}\) bis \(\pu{1 MeV}\) auf.
Dabei stößt ein Photon mit einem (quasi--freien) Elektron, wodurch es
einen Impuls auf das Elektron überträgt und gestreut wird.

Die Änderung der Wellenlänge \(\lambda\) des Photons hängt vom
Streuwinkel \(\theta\) ab. Der Wirkungsquerschnitt
\(\sigma_\mathrm{Co}\) ist von der Ordnungszahl \(Z\) des Absorbers und
der Energie \(E_\gamma\) der Strahlung abhängig. \([7]\)

\[
\begin{eqnarray}
    \Delta \lambda &=& \frac{h}{m_e c} (1 - \cos(\theta)) \\
    \sigma_\mathrm{Co} &\propto & \frac{Z}{E_\gamma}
\end{eqnarray}
\]

\hypertarget{paarbildung}{%
\subsubsection{Paarbildung}\label{paarbildung}}

Bei der Paarbildung entsteht aus einem Photon mit einer hohen Energie
ein Teilchen--Antiteilchen--Paar, dieser Effekt dominiert bei Energien
von etwa \(\pu{5 MeV}\). Die Umkehrung dieses Effektes nennt man
Paarvernichtung.

Im Coulombfeld eines Atomkernes kann ein Photon zu einem Positron
\(e^+\) und einem Elektron \(e^-\) zerstrahlen, wenn die Energie des
Photons mindestens \(\pu{1.022 MeV}\) beträgt. Diese Energie entspricht
der Gesamtmasse beider Teilchen.

Der Wirkungsquerschnitt \(\sigma_\mathrm{Paar}\) kann wie folgt
beschrieben werden.

\[
\begin{eqnarray}
    \sigma_\mathrm{Paar}
        &=& 4\alpha r_e^2 Z^2
            \left(\frac{7}{9} \ln(2\epsilon) - \frac{109}{54} \right) \\
    \sigma_\mathrm{Paar}
        &\propto& Z^2 \ln(E_\gamma)
\end{eqnarray}
\]

\hypertarget{szintillatoren-und-photomultiplier}{%
\subsection{Szintillatoren und Photomultiplier}\label{szintillatoren-und-photomultiplier}}

\clearpage
\hypertarget{durchfuxfchrung}{%
\section{Durchführung}\label{durchfuxfchrung}}

\clearpage
\hypertarget{auswertung}{%
\section{Auswertung}\label{auswertung}}

\clearpage
\hypertarget{fazit}{%
\section{Fazit}\label{fazit}}

\clearpage
\hypertarget{literatur}{%
\section{Literatur}\label{literatur}}

\begin{enumerate}
\def\labelenumi{\arabic{enumi}.}
\tightlist
\item
  Universität zu Köln, ``B3.4: Positronen--Emissions--Tomografie'',
  Januar 2021, Online verfügbar unter
  \url{https://www.ikp.uni--koeln.de/fileadmin/data/praktikum/B3.4\_PET\_de.pdf}
\item
  ``Chart of Nuclides'', National Nuclear Data Center,
  \url{https://www.nndc.bnl.gov/nudat3}, Abruf am 28.03.2024
\item
  ``Positronen Emissions Tomographie'', Deutsche Gesellschaft für
  Nuklearmedizin e.V., Online verfügbar unter
  \url{http://www.nuklearmedizin.de/docs/pet_bro_06.pdf}, Abruf am
  03.04.2024
\item
  W. Demtröder, ``Experimentalphysik 4: Kern-, Teilchen- und
  Astrophysik'', Springer--Spektrum--Verlag, 2017, DOI:
  \href{https://link.springer.com/book/10.1007/978-3-662-52884-6}{10.1007/978-3-662-52884-6}
\item
  Lexikon der Physik, ``Paarvernichtung'',
  \url{https://www.spektrum.de/lexikon/physik/paarvernichtung/10838},
  Abruf am 04.04.2024
\item
  Wikipedia, ``Annihilation'',
  \url{https://de.wikipedia.org/wiki/Annihilation}, Abruf am 04.04.2024
\item
  K. Bethge, Kernphysik: Eine Einführung, 3. aktualisierte und
  erweiterte Auflage, Springer--Verlag, 2008, DOI:
  \href{https://doi.org/10.1007/978-3-540-74567-9}{10.1007/978-3-540-74567-9}
\item
  Empfehlung der Strahlenschutzkommission, ``Orientierungshilfe für
  radiologische und nuklearmedizinische Untersuchungen'', In: ``Berichte
  der Strahlenschutzkommission (SSK) des Bundesministeriums für Umwelt,
  Naturschutz und Reaktorsicherheit: Heft~51'', 2006, ISBN~3873441306, \url{https://campus--nes.de/fileadmin/user\_upload/Orientierungshilfe\_SSK.pdf}
\end{enumerate}

\end{document}
